% =============================================================================
% Preámbulo compartido. Sigue el estilo de project_document/ para que los dos
% documentos del repositorio se lean como uno.
% =============================================================================

\documentclass[11pt,a4paper]{article}

\usepackage[utf8]{inputenc}
\usepackage[T1]{fontenc}
\usepackage[spanish,es-noquoting]{babel}
\usepackage[margin=2.4cm]{geometry}
\usepackage{microtype}
\usepackage{parskip}
\usepackage{enumitem}
\usepackage{xcolor}
\usepackage{booktabs}
\usepackage{longtable}
\usepackage{tabularx}
\usepackage{xltabular}
\usepackage{array}
\usepackage{amsmath}
\usepackage{amssymb}
\usepackage{titlesec}
\usepackage{fancyhdr}
\usepackage{caption}
\usepackage[hidelinks]{hyperref}

% --- color -------------------------------------------------------------------
\definecolor{acc}{HTML}{146B57}
\definecolor{dim}{HTML}{5C6B65}
\definecolor{rule}{HTML}{D8DEDA}
\definecolor{aviso}{HTML}{8A6410}

\hypersetup{
  unicode=true,
  colorlinks=true,
  linkcolor=acc,
  urlcolor=acc,
  citecolor=acc,
  pdftitle={PROTEA. Diseño de la campaña de experimentación},
  pdfauthor={Francisco Miguel Pérez Canales}
}

% --- títulos -----------------------------------------------------------------
\titleformat{\section}{\Large\bfseries\color{acc}}{\thesection}{0.8em}{}
\titleformat{\subsection}{\large\bfseries}{\thesubsection}{0.8em}{}
\titleformat{\subsubsection}{\normalsize\bfseries}{\thesubsubsection}{0.7em}{}
\titlespacing*{\section}{0pt}{2.2ex plus 1ex minus .2ex}{1.2ex plus .2ex}

% --- cabeceras ---------------------------------------------------------------
\pagestyle{fancy}
\fancyhf{}
\fancyhead[L]{\footnotesize\color{dim} PROTEA. Diseño de la campaña}
\fancyhead[R]{\footnotesize\color{dim}\thepage}
\renewcommand{\headrulewidth}{0.4pt}
\renewcommand{\headrule}{\hbox to\headwidth{\color{rule}\leaders\hrule height \headrulewidth\hfill}}

% --- tablas ------------------------------------------------------------------
\captionsetup{font=small,labelfont={bf,color=acc},skip=6pt}
\renewcommand{\arraystretch}{1.18}
\newcolumntype{L}[1]{>{\raggedright\arraybackslash}p{#1}}
\newcolumntype{C}[1]{>{\centering\arraybackslash}p{#1}}
\newcolumntype{R}[1]{>{\raggedleft\arraybackslash}p{#1}}

% --- el entorno de definiciones ---------------------------------------------
% Toda la sección 2 pasa por aquí. Un término definido aquí es el único
% significado que tiene en el resto del documento.
\newcounter{defn}
\makeatletter
% \term{Nombre}{cuerpo}  o  \term[etiqueta]{Nombre}{cuerpo}
% La etiqueta permite citar una definición con \ref, que imprime su número.
\newcommand{\term}[3][]{%
  \refstepcounter{defn}%
  \edef\pr@tmp{#1}%
  \ifx\pr@tmp\@empty\else\label{#1}\fi
  \par\vspace{0.65em}%
  \noindent\textbf{D\thedefn. #2.}\enspace #3\par\vspace{0.15em}}
\makeatother

% Un término del código. Se deja en inglés a propósito, véase la sección 2.0.
\newcommand{\cod}[1]{\texttt{\small #1}}

% Una cifra medida, con su procedencia.
\newcommand{\medido}[1]{\textcolor{dim}{\footnotesize (medido: #1)}}

% Un aviso que exige decision del investigador.
\newenvironment{aviso}[1]{%
  \par\vspace{0.9em}\noindent
  \begingroup\color{aviso}\rule{\linewidth}{1pt}\endgroup\par\vspace{0.35em}%
  \noindent\textbf{\color{aviso}#1}\par\vspace{0.25em}\noindent\ignorespaces}%
 {\par\vspace{0.35em}\noindent
  \begingroup\color{aviso}\rule{\linewidth}{1pt}\endgroup\par\vspace{0.9em}}
